% !TeX root = ./thuthesis-example.tex

% 论文基本信息配置（长沙理工大学硕士学位论文）

\thusetup{
  %******************************
  % 注意：
  %   1. 配置里面不要出现空行
  %   2. 不需要的配置信息可以删除
  %   3. 建议先阅读文档中所有关于选项的说明
  %******************************
  %
  % 输出格式
  output = electronic,
  %
  % 标题
  title  = {自动驾驶汽车自监督点云预测算法改善及应用研究},
  title* = {Research on the Improvement and Application of Self-Supervised Point Cloud Prediction Algorithms for Autonomous Vehicles},
  %
  % 学位类别
  degree-category  = {机械硕士},
  degree-category* = {Master of Engineering},
  %
  % 培养单位
  department = {汽车与机械工程学院},
  %
  % 学科
  discipline  = {机械工程},
  discipline* = {Mechanical Engineering},
  %
  % 姓名
  author  = {冯荣英},
  author* = {Feng Rongying},
  %
  % 指导教师
  supervisor  = {杜荣华, 教授},
  supervisor* = {Professor Du Ronghua},
  %
  % 副指导教师
  associate-supervisor  = {高凯, 副教授},
  associate-supervisor* = {Associate Professor Gao Kai},
  %
  % 日期
  date = {2024-04},
  %
  % 不生成书脊
  include-spine = false,
}

% 载入所需的宏包

% 定理类环境宏包
\usepackage{amsthm}

\thusetup{
  math-font  = xits,
}

% 表格加脚注
\usepackage{threeparttable}

% 表格中支持跨行
\usepackage{multirow}

% 跨页表格
\usepackage{longtable}

% 算法
\usepackage{algorithm}
\usepackage{algorithmic}

% 量和单位
\usepackage{siunitx}

% 参考文献使用 BibTeX + natbib 宏包（顺序编码制）
\usepackage[sort]{natbib}
\bibliographystyle{thuthesis-numeric}

% 定义所有的图片文件在 figures 子目录下
\graphicspath{{figures/}}

% 数学命令
\makeatletter
\newcommand\dif{%
  \mathop{}\!%
  \ifthu@math@style@TeX
    d%
  \else
    \mathrm{d}%
  \fi
}
\makeatother

% 载入长沙理工大学格式覆盖
\usepackage{csust-format}

% hyperref 宏包在最后调用
\usepackage{hyperref}
